\documentclass{article}
\usepackage{/globals/de}
\usepackage{/globals/anki}
\ankiprefix{wirtschaftliche-systeme}
\ankitopic{Wirtschaftliche Systeme}
\ankishow
\begin{document}
\section{Wirtschaftliche Systeme} 
Es gibt viele Anwendungssysteme mit spezifischen wirtschaftlichen Funktionen, für fast jeden wirtschaftlich relevanten Schritt des Vertriebs eines Produktes oder einer Dienstleistung.

Mit Hilfe von Computern können Produkte entworfen werden, Prozesse und Fabriken geplant werden, \etc. Es können Daten erfasst werden, der Fortschritt in der Produktion erfasst werden, Qualität gesichert werden, \etc.

\subsection{Bedarfsplanung}
Die Bedarfsplanung (\emph{MRP}, emph{Manufacturing Resource Plannig}) kann in zwei Phasen unterteilt werden. Zuerst wird bei der MRP II, der Primärbedarfsplanung, grob betrachtet wie viel Produziert werden kann bevor bei der MRP I, der Materialbedarfsplanung, betrachtet wird wie viele einzelne Materialien und Bestandteile dafür notwendig sind.
\ankinote{bedarfsplanung}{Bedarfsplanung}{MRP II (Primär), grob; MRP I (Material), Material}

Probleme bei der Fertigung können basierend auf der Losgröße antizipiert werden; ist sie zu hoch müssen Produkte lange gelagert werden, ist sie zu niedrig wird Kapazität und Zeit an Engpässen verschwendet.
\ankinote{losgröße}{Losgröße}{Kopien einer genauen Variation; klein -> verschwendete Kapazitäten, Wartung, groß -> Lagerung}

\subsection{Durchlaufterminierung}
Es wird bestimmt wann ein Arbeitsvorgang gestartet werden soll.
\begin{description}
 \item[Vorwärtsterminierung] Sobal der Auftrag erhalten wird.
 \item[Rückwärtsterminierung] So dass der Arbeitsvorgang genau dann fertig ist wenn das Produkt geliefert werden soll.
\end{description}
\ankinote{durchlaufterminierung}{Durchlaufterminierung}{Vorwärtsterminierung, sobald Auftrag erhalten; Rückwärtsterminierung, dass es zum Liefertermin fertig ist}
Aufträge und darin enthaltene Arbeitsänge können Zeitstrahlähnlich dargestellt werden.

Es sollte so geplant werden, dass Kapazitäten ausgeglichen sind.

\subsection{Verfügbarkeitsprüfung}
Es kann sichergestellt werden, dass die benötigten Ressourcen vorhanden sind; dass alle Materialien, Maschinen, Unterlagen, \etc vorhanden sind und funktionieren.

\subsection{Werkstattsteuerung}
Aufträge können anhand von unterschiedlichen Kriterien priorisiert werden.

\end{document}